% =====================================================================
\chapter{PROOF OF THE SUBSTITUTE LEMMA}
\label{app:proofs}
% =====================================================================

This appendix proves Lemma~\ref{lem:substitute}, on which the safety of the
local frontier (Theorem~\ref{thm:frontier}) rests: for every route $r$
outside the frontier and every admissible report $b$, there is a substitute
$u$ in the frontier (or the idle route) with $S_u\subseteq S_r$ and
$c_u(b)+q(S_r\setminus S_u)\le c_r(b)$.

\begin{proof}
Fix $r\notin\Kstar_i$ and $b\in\Theta$. On $\bar D=D(r)\cup\{r\}$ define
$\Phi(\rho)=c_\rho(b)+q(S_r\setminus S_\rho)$, the cost of serving $S_r$ with
route $\rho$ and sending the rest to FD. Since $\mu_r(b)\le 0$,
$m:=\min_{\bar D}\Phi\le\Phi(r)=c_r(b)$. Let $M=\arg\min_{\bar D}\Phi$ and
$k=\min_{\rho\in M}|S_\rho|$, the smallest bundle among the cheapest options.

\emph{Case $k=0$.} Then $0_i\in M$ and $u=0_i$ satisfies
$c_u(b)+q(S_r)=m\le c_r(b)$.

\emph{Case $k\ge 1$.} Fix $\rho_0\in M$ with $|S_{\rho_0}|=k$, let
$T=S_{\rho_0}$ and $M_T=\{\rho\in R_i:S_\rho=T,\;c_\rho(b)=c_{\rho_0}(b)\}$;
every $\rho\in M_T$ has $\Phi(\rho)=m$.

\emph{Claim.} For every $\rho\in M_T$ and every $\rho''\in D(\rho)\setminus M_T$,
$c_{\rho''}(b)+q(T\setminus S_{\rho''})>c_\rho(b)$.
If $S_{\rho''}\subsetneq T$, then $\rho''\in\bar D$ and $|S_{\rho''}|<k$, so
$\rho''\notin M$ and $\Phi(\rho'')>m=\Phi(\rho)$; since
$S_r\setminus S_{\rho''}=(T\setminus S_{\rho''})\cup(S_r\setminus T)$
disjointly, $\Phi(\rho'')-\Phi(\rho)=c_{\rho''}(b)+q(T\setminus S_{\rho''})-c_\rho(b)>0$.
If $S_{\rho''}=T$ and $\rho''\notin M_T$, then $\Phi(\rho'')\ge m$ gives
$c_{\rho''}(b)\ge c_\rho(b)$, and equality would place $\rho''$ in $M_T$.

The claim consists of finitely many strict inequalities between affine
functions of $b$, so there is $\eta>0$ such that they hold for every
$b'\in\Theta$ with $|b'-b|<\eta$. All routes of $M_T$ have bundle $T$ and the
same cost at $b$; by Assumption~\ref{ass:nodup} their slopes $W$ are
pairwise distinct. If $b<\uth$, let $u=\arg\min_{\rho\in M_T}W_\rho$ and choose
$b'\in(b,\min\{b+\eta,\uth\})$: then $c_u(b')<c_\rho(b')$ for all
$\rho\in M_T\setminus\{u\}$, and by the claim at $b'$,
$c_u(b')<c_{\rho''}(b')+q(T\setminus S_{\rho''})$ for all other
$\rho''\in D(u)$. Hence $\mu_u(b')>0$ and $u\in\Kstar_i$. If $b=\uth$, take
$u=\arg\max_{\rho\in M_T}W_\rho$ and $b'<b$ close to $b$; the same argument
applies. In both cases $u\ne r$, since otherwise $\mu_r(b')>0$ would
contradict $r\notin\Kstar_i$. Moreover $S_u=T\subseteq S_r$ and
$c_u(b)+q(S_r\setminus T)=m\le c_r(b)$.
\end{proof}
